\documentclass[a5paper,10pt]{article}

% https://github.com/InDevRus/filippov-solutions

% Нумерация страниц
\pagenumbering{gobble}

% Настройка полей
\usepackage{geometry}
\geometry{left=1cm}
\geometry{right=1cm}
\geometry{top=1cm}
\geometry{bottom=1.5cm}

% Вставка таблицы
\usepackage{array}
\usepackage{tabularx}

% Инструменты выравнивания формул
\usepackage{amsmath}
\usepackage{mathtools}

% Диакритика в математике
\usepackage{amsfonts}

% Вычисления размеров на ходу
\usepackage{calc}

% Удобные записи производных
\usepackage[italic=true]{derivative}

% Настройки сносок
\usepackage{enotez}

% Длинные знаки векторов
\usepackage{anyfontsize}
\usepackage[b]{esvect}

% Вставка картинок
\usepackage{graphicx}

% Вставка красной строки
\usepackage{indentfirst}
\setlength{\parindent}{1em}

% Условия в командах
\usepackage{xifthen}

% Луа-скрипты
\usepackage{luacode}

% Настройка команд отдельно в математическом и текстовом режимах
\usepackage{mathcommand}

% Для повторения LaTeX кода
\usepackage{multido}

% Конфигурация отступов формул
\delimitershortfall-1sp
\usepackage{mleftright}
\mleftright

% Растягивание объектов
\usepackage{relsize}

% Обтекание текстом
\usepackage{wrapfig}
\usepackage{subcaption}
\usepackage{floatflt}

% Поддержка ввода кириллицы
\usepackage[english, russian]{babel}

% Настройка корневой позиции для компилятора
\usepackage{import}
\providecommand*{\ProjectRootPath}{..}

\import{\ProjectRootPath/preambles/macros/}{theorems}

\import{\ProjectRootPath/preambles/fonts}{font_setup}

% Знак градуса
\usepackage{siunitx}

\import{\ProjectRootPath/preambles/macros/}{operators}
\import{\ProjectRootPath/preambles/macros/}{redefinitions}

\import{\ProjectRootPath/preambles/fonts}{letters_setup}
\import{\ProjectRootPath/preambles/fonts/adjustments}{custom_superscripts}

\import{\ProjectRootPath/preambles/custom}{formatting}
\import{\ProjectRootPath/preambles/custom}{frames}
\import{\ProjectRootPath/preambles/custom}{list_setup}

% TODO: перевести команды на современный стиль объявления

\begin{document}

$\dfrac {y\`} {\sqrt[3]{\pow{y}{2} + 1}} 
= \dfrac {1} {\sqrt[3]{\pow{x}{4} + 1}}$.
Общим решением будет набор функций

$$\tintlim{0}{y\br{x}} 
    {\dfrac {du} {\sqrt[3]{\pow{u}{2} + 1}}}
= \tintlim{0}{x} {\dfrac {\di t} {\sqrt[3]{\pow{t}{4} + 1}}} + C.$$

Интеграл 
$a = \tintlim{0}{+\infty} 
    {\dfrac {\di t} {\sqrt[3]{\pow{t}{4} + 1}}}$
сходится, так как подынтегральная функция знакопостоянна, и
$\dfrac {1} {\sqrt[3]{\pow{t}{4} + 1}} 
\sim \tpow{t}{\elevate[0.4em]{-\frac {4}{3}}}$
при $t \to +\infty$, а интеграл 
$\tintlim{1}{+\infty} {\tpow{t}{\elevate[0.4em]{-\frac {4}{3}}} \di t}$ сходится.
Точное значение этого интеграла равно
$a = \dfrac{1}{4}\Beta\br{\dfrac {1}{4}; \dfrac {1}{12}}$,
где $\Beta\br{x; y}$ -- бета-функция Эйлера.

Далее обозначим 
$g\br{\y\br{x}} 
= \tintlim{0}{y\br{x}} {\dfrac {du} {\sqrt[3]{\pow{u}{2} + 1}}}$.
Эта функция нечётная, так как подынтегральная функция чётная;
возрастает при $y > 0$, поскольку подынтегральная функция 
положительна; при $y\to +\infty$ бесконечно большая, 
поскольку интеграл 
$\tintlim{0}{+\infty} {\dfrac {du} {\sqrt[3]{\pow{u}{2} + 1}}}$
расходится. Действительно, при $u \to +\infty$ имеем 
$\dfrac {1} {\sqrt[3]{\pow{u}{2} + 1}} \sim \tpow{u}{\elevate[0.4em]{-\frac {2}{3}}}}$,
а соответствующий несобственный интеграл с такой
подынтегральной функцией расходится.

Из всего этого следует, что $g\br{\y}$ --
строго монотонно возрастающая непрерывная функция, 
принимающая в качестве значений все действительные числа. 
Значит существует обратная к ней функция 
$\pow{g}{\inv}\br{\y}$,
также сторого монотонно возрастающая, непрерывная 
и определённая для всех действительных $y$.

При этом $g\br{\y} = g\br{\y\br{x}}$ и
$\tlim{x \to \pm\infty} g\br{\y\br{x}}
= \tlim{x \to \pm\infty} g\br{\y\br{x}}
= \tlim{x \to \pm\infty}\ \br{\tintlim{0}{x} 
    {\dfrac {\di t} {\sqrt[3]{\pow{t}{4} + 1}}} + C}
= \linebreak
= \pm a + C$. Из-за непрерывности 
$g\br{\sub{y}{+\infty}} = a + C$ для некоторого 
$\sub{y}{+\infty}$, и $g\br{\sub{y}{-\infty}} = -a + C$ 
для некоторого $\sub{y}{-\infty}$. 
Следовательно, у всякого частного решения $y\br{x}$ 
есть горизонтальные асимптоты, равные 
$\pow{g}{\inv}\br{-a + C}$ и $\pow{g}{\inv}\br{a + C}$
при $x \to -\infty$ и $x \to +\infty$ соответственно.

\end{document}
